\documentclass[11pt,a4paper]{article}
\usepackage[utf8]{inputenc}
\usepackage[spanish]{babel}
\usepackage{amsmath}
\usepackage{amsfonts}
\usepackage{amssymb}
\usepackage{geometry}
\usepackage{hyperref}
\usepackage{booktabs}

\geometry{left=2.5cm, right=2.5cm, top=2.5cm, bottom=2.5cm}

\title{\textbf{Desarrollo Matemático Completo del Modelo de\\[4pt] Difusión-Advección-Segregación en Dunas Bidispersas}}
\author{Felipe Espinoza}
\date{\today}

\begin{document}

\maketitle

\tableofcontents
\newpage

% ======================================================================
%  PARTE I: GEOMETRÍA Y CINEMÁTICA
% ======================================================================

\section{El Problema Físico}
El modelo describe la evolución temporal y espacial de la concentración volumétrica de sedimentos finos, $\phi_s(x,z,t) \in [0,1]$, en el interior de una duna hidráulica bidispersa (compuesta por dos tamaños de partículas: finas de diámetro $d_s$ y gruesas de diámetro $d_l > d_s$). La dinámica del sistema está gobernada por tres mecanismos de transporte:
\begin{enumerate}
    \item \textbf{Advección:} Arrastre pasivo de las partículas por el campo de velocidades del fluido $(u, w)$.
    \item \textbf{Segregación gravitacional ($F_{seg}$):} Decantación de las partículas finas hacia el fondo del lecho activada por cizalle y gravedad.
    \item \textbf{Difusión cinemática ($D_{sl}$):} Remezcla caótica de partículas por colisiones inter-granulares.
\end{enumerate}

Se define la relación de tamaños como:
\begin{equation}
    R = \frac{d_l}{d_s}
\end{equation}

\section{Marco de Referencia Comóvil}
La duna migra a velocidad constante $c_{\text{mig}}$ en el marco de laboratorio. Para evitar una malla que se desplace en cada paso de tiempo, se trabaja en un \textbf{marco comóvil} donde la duna es estacionaria.

La relación entre ambos marcos es:
\begin{equation}
    x_{\text{lab}} = x' + c_{\text{mig}} \, t
\end{equation}
donde $x'$ es la coordenada en el marco de la duna. Todas las velocidades del fluido se corrigen restando la velocidad de migración: $u_{\text{duna}} = u_{\text{lab}} - c_{\text{mig}}$.

\section{Topografía de la Duna: $h(x')$}
La topografía del lecho se define como una función lineal por tramos en el marco de la duna. La duna tiene longitud $L_{\text{dune}}$, altura $H_d$ sobre un lecho base de altura $H_{\text{base}}$, y la cresta está ubicada en $x' = x_{\text{crest}}$. La longitud de la cara de sotavento es $L_{\text{lee}} = L_{\text{dune}} - x_{\text{crest}}$.

\begin{equation}
    h(x') = \begin{cases}
        H_{\text{base}} + H_d \displaystyle\frac{x'}{x_{\text{crest}}}, & 0 \le x' \le x_{\text{crest}} \quad \text{(Barlovento)} \\[8pt]
        H_{\text{base}} + H_d \left( 1 - \displaystyle\frac{x' - x_{\text{crest}}}{L_{\text{lee}}} \right), & x_{\text{crest}} < x' \le L_{\text{dune}} \quad \text{(Sotavento)}
    \end{cases}
\end{equation}

La pendiente topográfica $\frac{dh}{dx'}$ es constante por tramos:
\begin{equation}
    \frac{dh}{dx'} = \begin{cases}
        +\displaystyle\frac{H_d}{x_{\text{crest}}} > 0, & \text{Barlovento (pendiente suave ascendente)} \\[8pt]
        -\displaystyle\frac{H_d}{L_{\text{lee}}} < 0, & \text{Sotavento (pendiente pronunciada descendente)}
    \end{cases}
\end{equation}

\section{Función de Corriente y Campo de Velocidades}
Para garantizar que el campo de velocidades sea \textbf{estrictamente incompresible} ($\nabla \cdot \vec{u} = 0$), no se parametrizan $u$ y $w$ de forma independiente. En su lugar, se construye una \textbf{función de corriente} $\psi(x', z, t)$ a partir de la cual se derivan las velocidades.

\subsection{Definición de la Función de Corriente}
La función de corriente en el marco de laboratorio se define como:
\begin{equation}
    \psi_{\text{lab}}(x', z, t) = U_0 \, H_{\text{base}} \left( \frac{z}{h(x')} \right)^{m+1} g(x', t)
\end{equation}
donde:
\begin{itemize}
    \item $U_0$ es la velocidad de referencia de la corriente (escala del flujo).
    \item $H_{\text{base}}$ es la altura del lecho base.
    \item $m$ es el exponente del perfil de velocidades ($m > 0$; $m = 1$ corresponde a un perfil lineal, $m \to \infty$ a un perfil pistón).
    \item $g(x', t)$ es una función de modulación espacio-temporal que genera capas de deposición periódicas.
\end{itemize}

\subsection{Modulación Espacio-Temporal}
La función de modulación controla las variaciones periódicas del flujo en la zona de sotavento:
\begin{equation}
    g(x', t) = 1 + A_{\text{mod}} \sin\!\left(\frac{2\pi t}{T_{\text{period}}}\right) S_x(x')
\end{equation}
donde $A_{\text{mod}}$ es la amplitud de modulación, $T_{\text{period}}$ es el período, y $S_x(x')$ es un sigmoide suave que localiza la perturbación en la cara de sotavento:
\begin{equation}
    S_x(x') = \frac{1}{2} + \frac{1}{2} \tanh\!\left(\frac{x' - x_{\text{crest}}}{\delta_{\text{crest}}}\right)
\end{equation}
con derivada:
\begin{equation}
    \frac{dS_x}{dx'} = \frac{1}{2\delta_{\text{crest}}} \operatorname{sech}^2\!\left(\frac{x' - x_{\text{crest}}}{\delta_{\text{crest}}}\right)
\end{equation}
donde $\delta_{\text{crest}}$ es el ancho de la zona de transición suave en la cresta.

\subsection{Derivación de la Velocidad Horizontal $u$}
A partir de la definición de la función de corriente, la velocidad horizontal se obtiene como la derivada parcial respecto a $z$:
\begin{equation}
    u_{\text{lab}} = \frac{\partial \psi_{\text{lab}}}{\partial z}
\end{equation}
Derivando (5) respecto a $z$:
\begin{equation}
    u_{\text{lab}} = U_0 \, H_{\text{base}} \, (m+1) \frac{z^m}{h(x')^{m+1}} \, g(x', t) = U_0 \frac{H_{\text{base}}}{h(x')} (m+1) \eta^m \, g(x', t)
\end{equation}
donde se ha sustituido $z/h = \eta$.

En el \textbf{marco comóvil} de la duna, la velocidad horizontal corregida es:
\begin{equation}
    \boxed{u(x', \eta, t) = U_0 \frac{H_{\text{base}}}{h(x')} (m+1) \, \eta^m \, g(x', t) \; - \; c_{\text{mig}}}
\end{equation}

\subsection{Derivación de la Velocidad Vertical $w$}
La velocidad vertical se obtiene de la condición de incompresibilidad $\nabla \cdot \vec{u} = 0$, que equivale a:
\begin{equation}
    w = -\frac{\partial \psi_{\text{lab}}}{\partial x}
\end{equation}
Derivando (5) respecto a $x'$:
\begin{equation}
    w = -U_0 \, H_{\text{base}} \left[ \eta^{m+1} \frac{\partial g}{\partial x'} + g(x',t) (m+1) \eta^{m+1} \left( -\frac{1}{h} \frac{dh}{dx'} \right) \right]
\end{equation}

Simplificando:
\begin{equation}
    w = -U_0 \, H_{\text{base}} \, \eta^{m+1} \frac{\partial g}{\partial x'} + U_0 \frac{H_{\text{base}}}{h} (m+1) \eta^{m+1} g \frac{dh}{dx'}
\end{equation}

El segundo término se reconoce como $u_{\text{lab}} \cdot \eta \cdot \frac{dh}{dx'}$. Notar que la corrección comóvil del marco de la duna también aporta una contribución vertical:
\begin{equation}
    w_{\text{duna}} = w_{\text{lab}} + c_{\text{mig}} \eta \frac{dh}{dx'}
\end{equation}

\subsection{Velocidad Vertical en Coordenadas Sigma ($\omega_{\text{code}}$)}
En coordenadas sigma, la velocidad vertical ``de malla'' es la velocidad que cruza las superficies $\eta = \text{cte}$. Para el cálculo numérico se define:
\begin{equation}
    \omega_{\text{code}}(x', \eta, t) = c_{\text{mig}} \, \eta \frac{dh}{dx'} - U_0 \, H_{\text{base}} \, \eta^{m+1} \frac{\partial g}{\partial x'}
\end{equation}
donde:
\begin{equation}
    \frac{\partial g}{\partial x'} = A_{\text{mod}} \sin\!\left(\frac{2\pi t}{T_{\text{period}}}\right) \frac{dS_x}{dx'}
\end{equation}

\subsection{Verificación de Incompresibilidad}
El campo de velocidades es exactamente libre de divergencia en la formulación de volumen de control sigma:
\begin{equation}
    \frac{\partial (h \, u)}{\partial x'} + \frac{\partial \omega_{\text{code}}}{\partial \eta} = 0
\end{equation}
Esta identidad se satisface analíticamente por construcción, puesto que ambas velocidades se derivan de una misma función de corriente $\psi_{\text{lab}}$.

% ======================================================================
%  PARTE II: ECUACIÓN DE CONSERVACIÓN Y TRANSFORMACIÓN
% ======================================================================

\section{Ecuación Gobernante en Coordenadas Cartesianas}
Para un fluido incompresible, la ley de conservación de masa de la fase fina en coordenadas cartesianas $(x, z)$ es:
\begin{equation}
    \frac{\partial \phi_s}{\partial t} + \frac{\partial (u \phi_s)}{\partial x} + \frac{\partial (w \phi_s)}{\partial z} + \frac{\partial F_{seg}}{\partial z} = \frac{\partial}{\partial z} \left( D_{sl} \frac{\partial \phi_s}{\partial z} \right)
\end{equation}
donde:
\begin{itemize}
    \item El segundo y tercer término representan la \textbf{advección} (transporte pasivo por el fluido).
    \item $F_{seg} = -f_{sl} \, \phi_s(1-\phi_s)$ es el \textbf{flujo segregativo vertical} (decantación gravitacional). El factor $\phi_s(1-\phi_s)$ impide la segregación en zonas puras.
    \item El lado derecho es la \textbf{difusión cinemática} (remezcla por colisiones).
\end{itemize}

\section{Transformación a Coordenadas Adaptadas al Contorno ($\eta$)}
Para resolver la ecuación (19) en una topografía variable, transformamos las coordenadas cartesianas $(x, z, t)$ a coordenadas sigma $(x', \eta, t')$ que mapean el dominio a un rectángulo fijo.

\subsection{Definición de la Transformación}
Definimos el espesor total del lecho como $H(x,t) = h(x,t) - H_{\text{base}}$, y la transformación:
\begin{align}
    x' &= x \\
    \eta &= \frac{z - H_{\text{base}}}{H(x,t)} \\
    t' &= t
\end{align}

Bajo esta transformación: $\eta = 0$ corresponde a la base del lecho, $\eta = 1$ corresponde a la superficie de la duna, y el dominio físico variable se convierte en un rectángulo computacional fijo con $\eta \in [0, 1]$.

\subsection{Derivación de los Operadores Diferenciales}
Para transformar la EDP (19), expresamos las derivadas cartesianas en función de las coordenadas sigma mediante la regla de la cadena multivariable.

\subsubsection{Derivada Vertical}
\begin{equation}
    \frac{\partial}{\partial z} = \frac{\partial x'}{\partial z}\frac{\partial}{\partial x'} + \frac{\partial \eta}{\partial z}\frac{\partial}{\partial \eta} + \frac{\partial t'}{\partial z}\frac{\partial}{\partial t'}
\end{equation}

Dado que $x'$ y $t'$ no dependen de $z$: $\frac{\partial x'}{\partial z} = 0$ y $\frac{\partial t'}{\partial z} = 0$. Para la coordenada $\eta$:
\begin{equation}
    \frac{\partial \eta}{\partial z} = \frac{\partial}{\partial z}\left(\frac{z - H_{\text{base}}}{H}\right) = \frac{1}{H}
\end{equation}
Por lo tanto:
\begin{equation}
    \boxed{\frac{\partial}{\partial z} = \frac{1}{H}\frac{\partial}{\partial \eta}}
\end{equation}

\subsubsection{Derivada Horizontal}
\begin{equation}
    \frac{\partial}{\partial x} = \frac{\partial x'}{\partial x}\frac{\partial}{\partial x'} + \frac{\partial \eta}{\partial x}\frac{\partial}{\partial \eta} + \frac{\partial t'}{\partial x}\frac{\partial}{\partial t'}
\end{equation}

Aquí $\frac{\partial x'}{\partial x} = 1$ y $\frac{\partial t'}{\partial x} = 0$. Calculamos $\frac{\partial \eta}{\partial x}$ recordando que $z$ se mantiene fijo:
\begin{equation}
    \frac{\partial \eta}{\partial x} = \frac{\partial}{\partial x}\left(\frac{z - H_{\text{base}}}{H}\right) = (z - H_{\text{base}})\left(-\frac{1}{H^2}\frac{\partial H}{\partial x}\right) = -\frac{\eta}{H}\frac{\partial h}{\partial x}
\end{equation}
(ya que $H_{\text{base}}$ es constante, $\frac{\partial H}{\partial x} = \frac{\partial h}{\partial x}$). Sustituyendo:
\begin{equation}
    \boxed{\frac{\partial}{\partial x} = \frac{\partial}{\partial x'} - \frac{\eta}{H}\frac{\partial h}{\partial x}\frac{\partial}{\partial \eta}}
\end{equation}

\subsubsection{Derivada Temporal}
\begin{equation}
    \frac{\partial}{\partial t} = \frac{\partial x'}{\partial t}\frac{\partial}{\partial x'} + \frac{\partial \eta}{\partial t}\frac{\partial}{\partial \eta} + \frac{\partial t'}{\partial t}\frac{\partial}{\partial t'}
\end{equation}

Donde $\frac{\partial x'}{\partial t} = 0$ y $\frac{\partial t'}{\partial t} = 1$. Derivamos $\eta$ respecto al tiempo con $z$ fijo:
\begin{equation}
    \frac{\partial \eta}{\partial t} = (z - H_{\text{base}})\left(-\frac{1}{H^2}\frac{\partial H}{\partial t}\right) = -\frac{\eta}{H}\frac{\partial h}{\partial t}
\end{equation}
Sustituyendo:
\begin{equation}
    \boxed{\frac{\partial}{\partial t} = \frac{\partial}{\partial t'} - \frac{\eta}{H}\frac{\partial h}{\partial t}\frac{\partial}{\partial \eta}}
\end{equation}

\section{Sustitución Término a Término en la Ecuación de Conservación}
Aplicamos los operadores (25), (28) y (31) a cada término de la ecuación gobernante (19).

\subsection{Término Temporal}
\begin{equation}
    \frac{\partial \phi_s}{\partial t} = \frac{\partial \phi_s}{\partial t'} - \frac{\eta}{H}\frac{\partial h}{\partial t}\frac{\partial \phi_s}{\partial \eta}
\end{equation}

\subsection{Término de Advección Horizontal}
\begin{equation}
    \frac{\partial(u\phi_s)}{\partial x} = \frac{\partial(u\phi_s)}{\partial x'} - \frac{\eta}{H}\frac{\partial h}{\partial x}\frac{\partial(u\phi_s)}{\partial \eta}
\end{equation}

\subsection{Término de Advección Vertical}
\begin{equation}
    \frac{\partial(w\phi_s)}{\partial z} = \frac{1}{H}\frac{\partial(w\phi_s)}{\partial \eta}
\end{equation}

\subsection{Término de Segregación}
\begin{equation}
    \frac{\partial F_{seg}}{\partial z} = \frac{1}{H}\frac{\partial F_{seg}}{\partial \eta}
\end{equation}

\subsection{Término de Difusión}
Aplicando la derivada vertical dos veces:
\begin{equation}
    \frac{\partial}{\partial z}\left(D_{sl}\frac{\partial \phi_s}{\partial z}\right) = \frac{1}{H}\frac{\partial}{\partial \eta}\left(\frac{D_{sl}}{H}\frac{\partial \phi_s}{\partial \eta}\right)
\end{equation}

\subsection{Ecuación Completa Transformada}
Sumando los términos (32)--(35) e igualando al término difusivo (36):
\begin{multline}
    \left(\frac{\partial \phi_s}{\partial t'} - \frac{\eta}{H}\frac{\partial h}{\partial t}\frac{\partial \phi_s}{\partial \eta}\right) + \left(\frac{\partial(u\phi_s)}{\partial x'} - \frac{\eta}{H}\frac{\partial h}{\partial x}\frac{\partial(u\phi_s)}{\partial \eta}\right) \\
    + \frac{1}{H}\frac{\partial(w\phi_s)}{\partial \eta} + \frac{1}{H}\frac{\partial F_{seg}}{\partial \eta} = \frac{1}{H}\frac{\partial}{\partial \eta}\left(\frac{D_{sl}}{H}\frac{\partial \phi_s}{\partial \eta}\right)
\end{multline}

Multiplicamos toda la ecuación por $H(x,t)$ para eliminar los denominadores:
\begin{multline}
    H\frac{\partial \phi_s}{\partial t'} - \eta\frac{\partial h}{\partial t}\frac{\partial \phi_s}{\partial \eta} + H\frac{\partial(u\phi_s)}{\partial x'} - \eta\frac{\partial h}{\partial x}\frac{\partial(u\phi_s)}{\partial \eta} \\
    + \frac{\partial(w\phi_s)}{\partial \eta} + \frac{\partial F_{seg}}{\partial \eta} = \frac{\partial}{\partial \eta}\left(\frac{D_{sl}}{H}\frac{\partial \phi_s}{\partial \eta}\right)
\end{multline}

% ======================================================================
%  PARTE III: CAMBIO DE VARIABLE Q
% ======================================================================

\section{Introducción de la Variable de Masa Conservada $Q$}

\subsection{Motivación Física}
La masa total de finos en el sistema es:
\begin{equation}
    M_{\text{finos}} = \int_0^L \int_0^{h(x,t)} \phi_s \, dz \, dx = \int_0^L \int_0^1 \underbrace{h(x,t) \cdot \phi_s}_{Q} \, d\eta \, dx
\end{equation}

Al discretizar, cada celda computacional tiene dimensiones constantes $dx \times d\eta$. Si se resuelve directamente para $\phi_s$, la masa en cada celda es $\phi_s \cdot h(x,t) \cdot dx \cdot d\eta$, y al variar $h$ en el tiempo se generan errores de truncamiento graves. En cambio, si se define:
\begin{equation}
    Q(x', \eta, t') = H(x,t) \cdot \phi_s
\end{equation}
la masa de finos en cada celda es simplemente $Q \cdot dx \cdot d\eta$ (constante geométricamente). Al resolver para $Q$, se garantiza la conservación estricta de la masa.

\subsection{Relaciones de Producto para las Derivadas}
Como $H$ depende de $x'$ y $t'$ pero \textbf{no} de $\eta$, utilizamos la regla del producto:
\begin{align}
    \frac{\partial Q}{\partial t'} &= \frac{\partial(H\phi_s)}{\partial t'} = H\frac{\partial \phi_s}{\partial t'} + \phi_s\frac{\partial H}{\partial t'} \nonumber \\
    &\implies H\frac{\partial \phi_s}{\partial t'} = \frac{\partial Q}{\partial t'} - \phi_s\frac{\partial h}{\partial t}
\end{align}

\begin{align}
    \frac{\partial(uQ)}{\partial x'} &= \frac{\partial(Hu\phi_s)}{\partial x'} = H\frac{\partial(u\phi_s)}{\partial x'} + u\phi_s\frac{\partial H}{\partial x'} \nonumber \\
    &\implies H\frac{\partial(u\phi_s)}{\partial x'} = \frac{\partial(uQ)}{\partial x'} - u\phi_s\frac{\partial h}{\partial x}
\end{align}

\subsection{Sustitución en la Ecuación Transformada}
Sustituyendo (41) y (42) en la ecuación (38):
\begin{multline}
    \left(\frac{\partial Q}{\partial t'} - \phi_s\frac{\partial h}{\partial t}\right) - \eta\frac{\partial h}{\partial t}\frac{\partial \phi_s}{\partial \eta} + \left(\frac{\partial(uQ)}{\partial x'} - u\phi_s\frac{\partial h}{\partial x}\right) - \eta\frac{\partial h}{\partial x}\frac{\partial(u\phi_s)}{\partial \eta} \\
    + \frac{\partial(w\phi_s)}{\partial \eta} + \frac{\partial F_{seg}}{\partial \eta} = \frac{\partial}{\partial \eta}\left(\frac{D_{sl}}{H}\frac{\partial \phi_s}{\partial \eta}\right)
\end{multline}

\section{Agrupación en Flujos Conservativos}
Los términos que involucran derivadas de la topografía se agrupan de forma natural.

\subsection{Términos Temporales de la Topografía}
Tomamos los términos que dependen de $\frac{\partial h}{\partial t}$:
\begin{equation}
    -\left(\phi_s\frac{\partial h}{\partial t} + \eta\frac{\partial h}{\partial t}\frac{\partial \phi_s}{\partial \eta}\right)
\end{equation}
Dado que $\frac{\partial h}{\partial t}$ no depende de $\eta$, aplicamos la regla del producto inversa:
\begin{equation}
    \phi_s + \eta\frac{\partial \phi_s}{\partial \eta} = \frac{\partial(\eta \phi_s)}{\partial \eta}
\end{equation}
(verificación: $\frac{\partial(\eta \phi_s)}{\partial \eta} = \phi_s + \eta \frac{\partial \phi_s}{\partial \eta}$ \checkmark). Por lo tanto:
\begin{equation}
    -\frac{\partial h}{\partial t}\frac{\partial(\eta \phi_s)}{\partial \eta} = -\frac{\partial}{\partial \eta}\left(\eta \phi_s\frac{\partial h}{\partial t}\right)
\end{equation}

\subsection{Términos Espaciales de la Topografía}
Análogamente, los términos con $\frac{\partial h}{\partial x}$:
\begin{equation}
    -\left(u\phi_s\frac{\partial h}{\partial x} + \eta\frac{\partial h}{\partial x}\frac{\partial(u\phi_s)}{\partial \eta}\right)
\end{equation}
Factorizando $\frac{\partial h}{\partial x}$ (que no depende de $\eta$):
\begin{equation}
    u\phi_s + \eta\frac{\partial(u\phi_s)}{\partial \eta} = \frac{\partial(\eta u \phi_s)}{\partial \eta}
\end{equation}
Por lo tanto:
\begin{equation}
    -\frac{\partial h}{\partial x}\frac{\partial(\eta u \phi_s)}{\partial \eta} = -\frac{\partial}{\partial \eta}\left(\eta u \phi_s \frac{\partial h}{\partial x}\right)
\end{equation}

\subsection{Ecuación Consolidada}
Sustituyendo (46) y (49) en (43) y agrupando todas las derivadas $\partial/\partial\eta$ bajo un único operador:
\begin{equation}
    \frac{\partial Q}{\partial t'} + \frac{\partial(uQ)}{\partial x'} + \frac{\partial}{\partial \eta}\left[\underbrace{\left(w - \eta\frac{\partial h}{\partial t} - \eta u \frac{\partial h}{\partial x}\right)\phi_s}_{F_{adv}} + \underbrace{F_{seg}}_{F_{seg}} - \underbrace{\frac{D_{sl}}{H}\frac{\partial \phi_s}{\partial \eta}}_{F_{diff}}\right] = 0
\end{equation}

Esta es la \textbf{forma conservativa estricta} de la ecuación de transporte en coordenadas sigma. Garantiza la conservación global de masa de la fase fina ante cualquier deformación de la topografía.

% ======================================================================
%  PARTE IV: FLUJOS FÍSICOS — MODELO DE TREWHELA ET AL. (2021)
% ======================================================================

\section{Modelación de los Flujos Físicos: Trewhela et al.\ (2021)}
La ecuación (50) requiere cerrar las expresiones para los flujos segregativo y difusivo. Utilizamos el modelo de Trewhela et al.\ (2021), que los parametriza en función del cizalle del flujo, la composición local y la presión litostática.

\subsection{Tasa de Cizalle Local $\dot{\gamma}$}
La tasa de cizalle del fluido se calcula derivando la velocidad horizontal respecto a la coordenada vertical:
\begin{equation}
    \dot{\gamma}(\eta) = \left|\frac{\partial u}{\partial z}\right| = \frac{1}{H}\left|\frac{\partial u}{\partial \eta}\right|
\end{equation}
A partir de la ecuación (11):
\begin{equation}
    \dot{\gamma}(\eta) = \frac{U_0 \, H_{\text{base}}}{h^2(x')} \, m(m+1) \, \eta^{m-1} \, g(x', t)
\end{equation}

\subsection{Diámetro Medio Local $\bar{d}$}
El diámetro medio de la mezcla bidispersa en función de la concentración local:
\begin{equation}
    \bar{d}(\phi_s) = (1 - \phi_s) \, d_l + \phi_s \, d_s
\end{equation}

\subsection{Velocidad de Segregación $f_{sl}$}
La velocidad de segregación depende del cizalle, la composición y la presión litostática:
\begin{equation}
    f_{sl} = \frac{B \, \dot{\gamma} \, \bar{d}^2}{C \, \bar{d} + p(\eta)} \cdot F(R, \phi_s)
\end{equation}
donde:
\begin{itemize}
    \item $B = 0.3744$ y $C = 0.2712$ son constantes empíricas.
    \item $p(\eta) = \Phi_s \, h(x') \, (1 - \eta)$ es la presión litostática normalizada, donde $\Phi_s \approx 0.6$ es la fracción volumétrica de sólidos.
    \item $F(R, \phi_s)$ es la función de asimetría de segregación:
    \begin{equation}
        F(R, \phi_s) = (R - 1) + E \, (1 - \phi_s)(R - 1)^2
    \end{equation}
    con $E = 2.0957$. La dependencia en $(1 - \phi_s)$ refleja que la segregación es más intensa cuando hay más espacio intergranular (más partículas gruesas).
\end{itemize}

El flujo segregativo vertical completo es:
\begin{equation}
    F_{seg} = -f_{sl} \, \phi_s(1 - \phi_s)
\end{equation}

\subsection{Coeficiente de Difusión $D_{sl}$}
La difusión cinemática escala cuadráticamente con el tamaño medio y linealmente con el cizalle:
\begin{equation}
    D_{sl} = A \, \dot{\gamma} \, \bar{d}^2
\end{equation}
donde $A = 0.01$ es una constante adimensional. El flujo difusivo en la coordenada sigma es:
\begin{equation}
    F_{diff} = \frac{D_{sl}}{H}\frac{\partial \phi_s}{\partial \eta}
\end{equation}

\subsection{Tabla de Parámetros del Modelo de Trewhela}
\begin{table}[h!]
\centering
\caption{Constantes del modelo de Trewhela et al.\ (2021)}
\vspace{0.2cm}
\begin{tabular}{lccl}
\toprule
\textbf{Parámetro} & \textbf{Símbolo} & \textbf{Valor} & \textbf{Descripción} \\
\midrule
Constante de difusión    & $A$ & $0.01$   & Adimensional \\
Constante de segregación & $B$ & $0.3744$ & Adimensional \\
Constante de presión     & $C$ & $0.2712$ & Adimensional \\
Constante de asimetría   & $E$ & $2.0957$ & Adimensional \\
Fracción de sólidos      & $\Phi_s$ & $0.6$ & Compactación del lecho \\
\bottomrule
\end{tabular}
\end{table}

% ======================================================================
%  PARTE V: ECUACIÓN FINAL DEL SOLVER
% ======================================================================

\section{Ecuación Final Implementada en el Código}
Combinando la forma conservativa (50) con los flujos físicos del modelo de Trewhela, la ecuación que resuelve el código de simulación es:
\begin{equation}
    \boxed{\frac{\partial Q}{\partial t'} + \frac{\partial(uQ)}{\partial x'} + \frac{\partial F_\eta}{\partial \eta} = 0}
\end{equation}
donde el flujo vertical total numérico se descompone como:
\begin{equation}
    F_\eta = F_{adv, \eta} + F_{seg, \eta} - F_{diff, \eta}
\end{equation}

Con cada componente definida explícitamente como:
\begin{align}
    F_{adv,\eta} &= \left(w - \eta\frac{\partial h}{\partial t} - \eta u \frac{\partial h}{\partial x}\right)\phi_s = (H \, w_\eta) \, \phi_s \\
    F_{seg,\eta} &= -f_{sl} \, \phi_s(1 - \phi_s) \\
    F_{diff,\eta} &= \frac{D_{sl}}{H}\frac{\partial \phi_s}{\partial \eta}
\end{align}

La velocidad vertical en coordenadas sigma $w_\eta$ se define como:
\begin{equation}
    w_\eta = \frac{d\eta}{dt} = \frac{1}{H}\left(w - \eta\frac{\partial h}{\partial t} - \eta u \frac{\partial h}{\partial x}\right)
\end{equation}

% ======================================================================
%  PARTE VII: CONDICIONES DE BORDE
% ======================================================================

\section{Condiciones de Borde del Sistema}

\subsection{Base del Lecho ($\eta = 0$)}
Condición impermeable: flujo neto nulo a través de la base.
\begin{equation}
    F_\eta\big|_{\eta=0} = 0
\end{equation}

\subsection{Superficie de la Duna ($\eta = 1$)}
La superficie es una frontera abierta con balance de masa dinámico que depende del signo de la velocidad vertical:
\begin{itemize}
    \item \textbf{Erosión (barlovento, $\omega_{\text{code}} \ge 0$):} Las partículas salen de la superficie con la concentración local:
    \begin{equation}
        F_\eta\big|_{\eta=1} = w_\eta \, \phi_s(x', 1, t)
    \end{equation}
    \item \textbf{Deposición (sotavento, $\omega_{\text{code}} < 0$):} Las partículas entran con una concentración modulada:
    \begin{equation}
        \phi_{\text{inflow}}(x', t) = \text{clip}\Big(\lambda(t) \cdot P_{\text{spatial}}(x') \cdot P_{\text{temporal}}(t),\; 0,\; 1\Big)
    \end{equation}
\end{itemize}

La función $P_{\text{spatial}}$ controla la clasificación granulométrica espacial sobre la cara de sotavento:
\begin{equation}
    P_{\text{spatial}}(x') = \max\!\left(0.1,\; 1 + \alpha_{\text{spatial}}\left(0.5 - \frac{x' - x_{\text{crest}}}{L_{\text{lee}}}\right)\right)
\end{equation}

La función $P_{\text{temporal}}$ controla la clasificación temporal periódica (creación de capas):
\begin{equation}
    P_{\text{temporal}}(t) = 1 + A_P\sin\!\left(\frac{2\pi t}{T_{\text{period}}}\right)
\end{equation}

El multiplicador de masa $\lambda(t)$ se calcula para imponer un balance global estricto entre la masa erosionada (barlovento) y la masa depositada (sotavento) en cada instante:
\begin{equation}
    \lambda(t) = \frac{\displaystyle\int_{\text{stoss}} \omega_{\text{code}}(x') \, \phi_s(x', 1, t) \, dx'}{\displaystyle\int_{\text{lee}} |\omega_{\text{code}}(x')| \, P_{\text{spatial}}(x') \, P_{\text{temporal}}(t) \, dx'}
\end{equation}

\subsection{Condiciones Horizontales (Periódicas)}
El dominio es periódico en $x'$: el material que sale por la frontera derecha ($x' = L_{\text{dune}}$) reingresa por la izquierda ($x' = 0$), cerrando el ciclo de reciclaje de partículas.

% ======================================================================
%  PARTE VIII: MÉTODO NUMÉRICO
% ======================================================================

\section{Método Numérico}

\subsection{Integración Temporal: SSP-RK2}
Se utiliza un esquema Runge-Kutta de 2$^\circ$ orden de estabilidad fuerte (\emph{Strong Stability Preserving}, SSP-RK2) para evitar la dispersión numérica:
\begin{align}
    Q^* &= Q^n + \Delta t \, \mathcal{R}(Q^n, t^n) \\
    Q^{n+1} &= \frac{1}{2}Q^n + \frac{1}{2}\left(Q^* + \Delta t \, \mathcal{R}(Q^*, t^n + \Delta t)\right)
\end{align}
donde $\mathcal{R}$ denota el lado derecho completo de la ecuación.

\subsection{Advección Horizontal: Esquema Upwind}
El flujo horizontal se evalúa en las caras de las celdas usando un esquema \emph{upwind} de primer orden:
\begin{equation}
    F_{x, i+1/2} = \begin{cases}
        u_{i+1/2} \, Q_i, & \text{si } u_{i+1/2} \ge 0 \\
        u_{i+1/2} \, Q_{i+1}, & \text{si } u_{i+1/2} < 0
    \end{cases}
\end{equation}

\subsection{Condición CFL}
El paso de tiempo se restringe por la condición de Courant-Friedrichs-Lewy:
\begin{equation}
    \Delta t \le 0.8 \left(\frac{u_{\max}}{dx} + \frac{w_{\eta,\max} + f_{sl,\max}/h_{\min}}{d\eta}\right)^{-1}
\end{equation}

Si la difusión está activada, se impone adicionalmente:
\begin{equation}
    \Delta t \le 0.45 \, \frac{(d\eta \cdot h_{\min})^2}{D_{sl,\max}}
\end{equation}

\subsection{Corrección de Masa Global}
Al final de cada paso RK2, se aplica una corrección multiplicativa para eliminar errores de redondeo acumulados:
\begin{equation}
    Q^{n+1}_{\text{final}} = Q^{n+1} \cdot \frac{M_{\text{initial}}}{M_{\text{current}}}
\end{equation}
donde $M_{\text{initial}} = \sum_{i,j} Q_{i,j} \, dx \, d\eta$ es la masa total inicial del sistema.

\end{document}
